\chapter*{读码检查清单}
\addcontentsline{toc}{chapter}{读码检查清单}

\begin{description}
  \item[构建入口] 是否能从 CMake 文件看出每个可执行目标依赖哪些源码，哪些目标属于默认 CTest，哪些目标依赖外部环境。
  \item[资产边界] 模型文件、tokenizer、safetensors header 和外部下载模型是否都有校验或拒绝路径。
  \item[公共合同] 头文件是否把 shape、dtype、token、trace checkpoint、KV cache 地址和返回字段说清楚。
  \item[正确性路径] tiny MLP、reference decoder 和 GPT-2 tiny fixture 是否有可复查 golden。
  \item[性能口径] benchmark 是否同时报告 backend、available、耗时、误差和 checksum，避免只有速度没有正确性。
  \item[真实模型证据] GPT-2 smoke 和 SmolLM2 smoke 是否记录模型文件 hash、命令、输出和验证结论。
  \item[回归门禁] CTest 是否覆盖坏资产、坏 shape、tokenizer 合同、trace 关键字段、ledger 固定公式和 serving probe 拒绝路径。
\end{description}

如果某段新代码不能回答上面的检查项，它就还不能进入第三册主线工程。源码卷的作用是把这条标准显式化：代码不只是能编译，还要能解释、能测量、能拒绝坏输入、能被长期回归。
